\section{未来方向与开放问题}

\subsection{单 Agent vs. 多 Agent 架构}

Anthropic 指出，目前尚不清楚单一的通用编码 Agent 是否在所有 context 下都能表现最佳，还是通过多 Agent 架构可以实现更好的性能。可能的专业化 Agent 包括：

\begin{itemize}
    \item \textbf{Testing Agent}：专门负责编写和运行测试
    \item \textbf{QA Agent}：专门负责质量保证和 bug 检测
    \item \textbf{Code Cleanup Agent}：专门负责代码重构和清理
    \item \textbf{Architecture Agent}：专门负责系统设计和架构决策
\end{itemize}

\subsection{领域泛化}

当前的方案主要针对全栈 Web 应用开发进行了优化。Anthropic 认为这些经验可能适用于其他领域，包括但不限于：

\begin{itemize}
    \item 科学研究中的实验设计和数据分析
    \item 金融建模与量化分析
    \item 系统运维和基础设施管理
    \item 文档撰写和内容创作
\end{itemize}

\begin{knowledgebox}{从 Web 开发到通用场景}
Web 应用开发之所以成为首选实验场景，是因为它具备几个有利条件：功能可以被清晰地分解为独立单元，端到端测试可以通过浏览器自动化实现，且开发流程（编码-测试-提交）高度标准化。将这些实践泛化到其他领域时，关键挑战在于如何定义"功能"和"测试"——在科学研究中，一个"功能"可能是"验证一个假设"，而"测试"可能是"重现实验结果"。
\end{knowledgebox}

\subsection{本章小结}

长时间运行 Agent 的研究仍处于早期阶段。多 Agent 架构是否优于单 Agent，以及这些实践如何泛化到 Web 开发之外的领域，都是值得深入探索的方向。


